\subsection{Generated automata with disjunctive invariants}
The following two tracked inputs show how the export handles disjunctive
invariants in automata produced by the integrated tool. They illustrate the
verified post-processing and export layer; the upstream translation is not
part of the generic preservation claims in this paper.

\subsubsection{Nested upper-bounded Until}
For
\[
  \Diamond_{\le5}(p\,\mathsf U_{\le2}\,q),
\]
the optimized automaton has a disjunctive invariant at location~4,
$x_{u1}\le2\lor x_{u0}\le5$. Every edge entering that location resets
$x_{u1}$. If the entry value of $x_{u0}$ is $v$, the two invariant clauses
permit delays of $2$ and $5-v$. The export splits location~4 into a copy with
invariant $x_{u1}\le2$ and a copy with invariant $x_{u0}\le5$. Their entry
conditions are respectively $x_{u0}\ge3$ and $x_{u0}\le3$: each selects a
copy that admits a maximal delay. Both are enabled at the tie $x_{u0}=3$.
Incoming transitions, including the self-loop, are replicated to both
copies, and outgoing transitions are replicated from each.

Figure~\ref{fig:ex1-invariants} shows the generated automaton before and after
the split. For readability, the drawing omits guards implied by source
invariants and edges subsumed by another edge with the same target, event,
and resets; the exported automaton retains them. The output has 23 explicit
transitions with \texttt{-init} (25 without it). This example uses Spot
options \texttt{-B -D}; dead resets have already been removed.

\begin{figure*}[!tp]
\centering
\subfloat[Before invariant elimination.\label{fig:ex1-with-disj}]{
  \begin{tikzpicture}[x=1cm,y=1cm]
  \node[fmstate,minimum size=13mm] (s1) at (0,0) {$\begin{array}{c}1\\[-1mm]x_{u0}\le5\end{array}$};
  \node[fmstate,minimum size=18mm] (s4) at (4.4,3.0) {$\begin{array}{c}4\\[-1mm]x_{u1}\le2\;\lor\;x_{u0}\le5\end{array}$};
  \node[fmstate,minimum size=14mm] (s3) at (8.3,-0.5) {$\begin{array}{c}3\\[-1mm]x_{u1}\le2\end{array}$};
  \node[fmaccept] (s2) at (12.4,2.2) {$2$};
  \draw[fmedge] (-1.1,0) -- (s1);
  \draw[fmedge] (s1) edge[loop above,looseness=4] node[fmlabel] {$\mathit{other}/\mathsf{rst}(x_{u1})$} (s1);
  \draw[fmedge] (s1) to[bend left=8] node[fmlabel,pos=.45,above,sloped] {$p/\mathsf{rst}(x_{u1})$} (s4);
  \draw[fmedge] (s4) to[bend left=13] node[fmlabel,pos=.50,below,sloped] {$\begin{array}{c}x_{u0}\le5,\;\mathit{other}/\\\mathsf{rst}(x_{u1})\end{array}$} (s1);
  \draw[fmedge] (s1) to[bend right=8] node[fmlabel,pos=.63,below] {$p/\mathsf{rst}(x_{u1})$} (s3);
  \draw[fmedge] (s1) to[bend right=48] node[fmlabel,pos=.4,below,sloped] {$q$} (s2);
  \draw[fmedge] (s4) edge[loop above,looseness=5] node[fmlabel] {$\begin{array}{c}x_{u0}\le5\\p/\mathsf{rst}(x_{u1})\end{array}$} (s4);
  \draw[fmedge] (s4) to[bend left=21] node[fmlabel,pos=.54,above,sloped] {$x_{u0}\le5,\;q$} (s2);
  \draw[fmedge] (s4) to[bend left=8] node[fmlabel,pos=.56,below,sloped] {$x_{u1}\le2,\;q$} (s2);
  \draw[fmedge] (s4) to[bend left=5] node[fmlabel,pos=.53,above,sloped] {$\begin{array}{c}x_{u0}\le5\\p/\mathsf{rst}(x_{u1})\end{array}$} (s3);
  \draw[fmedge] (s4) to[bend right=32] node[fmlabel,pos=.5,below,sloped] {$x_{u1}\le2,\;p$} (s3);
  \draw[fmedge] (s3) edge[loop right,looseness=5] node[fmlabel,right] {$p$} (s3);
  \draw[fmedge] (s3) -- node[fmlabel,above,sloped] {$q$} (s2);
  \draw[fmedge] (s2) edge[loop right,looseness=5] node[fmlabel,right] {$p,\;q,\;\mathit{other}$} (s2);
\end{tikzpicture}
}

\medskip

\subfloat[After invariant elimination.\label{fig:ex1-no-disj}]{
  \begin{tikzpicture}[x=1cm,y=1cm]
  % Location 4 of the original automaton is split into 4 (Inv x_u1<=2) and 5 (Inv x_u0<=5).
  % Every entry into the split location resets x_u1, so copy 4 allows a delay of 2 and
  % copy 5 a delay of 5 - x_u0; the entry guards select the copy with the longer delay.
  \node[fmstate,minimum size=13mm] (s1) at (0,0) {$\begin{array}{c}1\\[-1mm]x_{u0}\le5\end{array}$};
  \node[fmstate,minimum size=14mm] (s4) at (4.6,2.5) {$\begin{array}{c}4\\[-1mm]x_{u1}\le2\end{array}$};
  \node[fmstate,minimum size=14mm] (s5) at (4.6,-2.5) {$\begin{array}{c}5\\[-1mm]x_{u0}\le5\end{array}$};
  \node[fmstate,minimum size=14mm] (s3) at (9.4,0) {$\begin{array}{c}3\\[-1mm]x_{u1}\le2\end{array}$};
  \node[fmaccept] (s2) at (13.2,0) {$2$};
  \draw[fmedge] (-1.2,0) -- (s1);
  % location 1
  \draw[fmedge] (s1) edge[loop below,looseness=4] node[fmlabel,below] {$\mathit{other}/\mathsf{rst}(x_{u1})$} (s1);
  \draw[fmedge] (s1) to[bend left=10] node[fmlabel,pos=.45,above,sloped] {$x_{u0}\ge3,\;p/\mathsf{rst}(x_{u1})$} (s4);
  \draw[fmedge] (s1) to[bend right=10] node[fmlabel,pos=.45,below,sloped] {$x_{u0}\le3,\;p/\mathsf{rst}(x_{u1})$} (s5);
  \draw[fmedge] (s1) -- node[fmlabel,pos=.9,below] {$p/\mathsf{rst}(x_{u1})$} (s3);
  \draw[fmedge] (s1) to[bend left=88] node[fmlabel,above] {$q$} (s2);
  % location 4 (Inv x_u1 <= 2)
  \draw[fmedge] (s4) to[bend left=10] node[fmlabel,pos=.55,below,sloped] {$x_{u0}\le5,\;\mathit{other}/\mathsf{rst}(x_{u1})$} (s1);
  \draw[fmedge] (s4) edge[loop above,looseness=4] node[fmlabel,above] {$3\le x_{u0}\le5,\;p/\mathsf{rst}(x_{u1})$} (s4);
  \draw[fmedge] (s4) to[bend left=14] node[fmlabel,right] {$\begin{array}{l}x_{u0}\le3,\\p/\mathsf{rst}(x_{u1})\end{array}$} (s5);
  \draw[fmedge] (s4) to[bend left=10] node[fmlabel,pos=.5,above,sloped] {$x_{u0}\le5,\;p/\mathsf{rst}(x_{u1})$} (s3);
  \draw[fmedge] (s4) to[bend right=12] node[fmlabel,pos=.55,below,sloped] {$p$} (s3);
  \draw[fmedge] (s4) to[bend left=28] node[fmlabel,pos=.6,above,sloped] {$q$} (s2);
  % location 5 (Inv x_u0 <= 5)
  \draw[fmedge] (s5) to[bend right=10] node[fmlabel,pos=.55,above,sloped] {$\mathit{other}/\mathsf{rst}(x_{u1})$} (s1);
  \draw[fmedge] (s5) edge[loop below,looseness=4] node[fmlabel,below] {$x_{u0}\le3,\;p/\mathsf{rst}(x_{u1})$} (s5);
  \draw[fmedge] (s5) to[bend left=14] node[fmlabel,left] {$\begin{array}{r}x_{u0}\ge3,\\p/\mathsf{rst}(x_{u1})\end{array}$} (s4);
  \draw[fmedge] (s5) to[bend right=10] node[fmlabel,pos=.5,below,sloped] {$p/\mathsf{rst}(x_{u1})$} (s3);
  \draw[fmedge] (s5) to[bend right=30] node[fmlabel,pos=.6,below,sloped] {$x_{u1}\le2,\;p$} (s3);
  \draw[fmedge] (s5) to[bend right=28] node[fmlabel,pos=.6,below,sloped] {$q$} (s2);
  % locations 3 and 2
  \draw[fmedge] (s3) edge[loop below,looseness=5] node[fmlabel,below] {$p$} (s3);
  \draw[fmedge] (s3) -- node[fmlabel,above] {$q$} (s2);
  \draw[fmedge] (s2) edge[loop right,looseness=5] node[fmlabel,right] {$p,\;q,\;\mathit{other}$} (s2);
\end{tikzpicture}
}
\caption{Generated automaton for $\Diamond_{\le5}(p\,\mathsf U_{\le2}\,q)$.
Location invariants appear inside the nodes; edge labels are ``guard, event /
resets.'' A comma-separated event list denotes separate transitions, and
$\mathit{other}$ denotes any event other than $p$ and $q$. The double circle
marks the accepting location. Both diagrams are after the initial
\texttt{\_init\_} transition at time zero. Location~5 is created by the
split. Guards implied by source invariants and subsumed edges are omitted from
the drawing.}
\label{fig:ex1-invariants}
\end{figure*}
At the timed-word level, the accepted behavior has a direct witness: either
$q$ occurs at a position $k$ with $t_k\le5$, or a sequence satisfying $p$
starts at some $j$ with $t_j\le5$ and reaches $q$ at $k$ with
$t_k-t_j\le2$. The split preserves these same choices at location~4 while
making each location invariant conjunctive.

\subsubsection{Timed Release obligation inside Until}
The second input is
\[
  \Diamond_{\le5}\bigl((\Box_{\le3}(p\lor q))\,\mathsf U_{\le2}\,q\bigr).
\]
Its optimized automaton has seven locations before export and contains a
disjunctive invariant (Fig.~\ref{fig:ex2-disj}). The three distinct primitive
timed subformulas use clocks $x_{u0}$, $x_{u1}$, and $x_{r2}$. Two syntactic
occurrences of the bounded Release share $x_{r2}$ because the clock key is
the primitive formula.

When a $q$ transition enters accepting location~5, the automaton permits only
$p$ or $q$ until $x_{r2}>3$, after which location~2 accepts every event. The
bounded Release obligation can therefore remain pending after the Until
witness. Two $q$ transitions testing $x_{r2}>3$ are unsatisfiable: at
location~4, entry resets both relevant clocks; at location~3, entry resets
$x_{r2}$ and preserves $x_{u1}$ with $x_{r2}\le x_{u1}\le2$. The verified
passes compare clock bounds with constants and do not track this clock-order
relation, so the transitions remain in the optimized drawing. They cannot be
taken. Export produces seven locations and 29 transitions.

\begin{figure*}[!t]
\centering
\resizebox{\textwidth}{!}{\begin{tikzpicture}[x=1cm,y=1cm]
  % Output of mtl2tba -init after optimization, before the export
  % (clocks printed x0, x1, x2 by the tool are x_u0, x_r2, x_u1).
  \node[fmstate,minimum size=13mm] (s1) at (0,0) {$\begin{array}{c}1\\[-1mm]x_{u0}\le5\end{array}$};
  \node[fmstate,minimum size=19mm] (s4) at (4.6,3.4) {$\begin{array}{c}4\\[-1mm]x_{u1}\le2\;\lor\;x_{u0}\le5\end{array}$};
  \node[fmstate,minimum size=14mm] (s3) at (8.4,-0.6) {$\begin{array}{c}3\\[-1mm]x_{u1}\le2\end{array}$};
  \node[fmaccept] (s5) at (11.2,3.4) {$5$};
  \node[fmaccept] (s2) at (13.6,0.4) {$2$};
  \draw[fmedge] (-1.1,0) -- (s1);
  % location 1
  \draw[fmedge] (s1) edge[loop below,looseness=4] node[fmlabel,below] {$p,\;\mathit{other}/\mathsf{rst}(x_{u1})$} (s1);
  \draw[fmedge] (s1) to[bend left=8] node[fmlabel,pos=.42,above,sloped] {$p/\mathsf{rst}(x_{r2}),\mathsf{rst}(x_{u1})$} (s4);
  \draw[fmedge] (s1) to[bend right=6] node[fmlabel,pos=.55,below] {$p/\mathsf{rst}(x_{r2}),\mathsf{rst}(x_{u1})$} (s3);
  \draw[fmedge] (s1) to[bend right=38] node[fmlabel,pos=.5,below] {$q$} (s2);
  % location 4
  \draw[fmedge] (s4) to[bend left=14] node[fmlabel,pos=.5,below,sloped] {$\begin{array}{c}x_{u0}\le5,\;p,\;\mathit{other}\\/\mathsf{rst}(x_{u1})\end{array}$} (s1);
  \draw[fmedge] (s4) edge[loop above,looseness=5] node[fmlabel] {$\begin{array}{c}x_{u0}\le5,\;p\\/\mathsf{rst}(x_{r2}),\mathsf{rst}(x_{u1})\end{array}$} (s4);
  \draw[fmedge] (s4) to[bend left=10] node[fmlabel,pos=.5,above,sloped] {$\begin{array}{c}x_{u0}\le5,\;p\\/\mathsf{rst}(x_{r2}),\mathsf{rst}(x_{u1})\end{array}$} (s3);
  \draw[fmedge] (s4) to[bend right=22] node[fmlabel,pos=.5,below,sloped] {$x_{u1}\le2,\;p/\mathsf{rst}(x_{r2})$} (s3);
  \draw[fmedge] (s4) -- node[fmlabel,above] {$x_{u1}\le2,\;q$} (s5);
  \draw[fmedge] (s4) to[bend left=6] node[fmlabel,pos=.62,above,sloped] {$x_{u0}\le5,\;q$} (s2);
  \draw[fmedge] (s4) to[bend right=4] node[fmlabel,pos=.62,below,sloped] {$x_{r2}>3\land x_{u1}\le2,\;q$} (s2);
  % location 3
  \draw[fmedge] (s3) edge[loop below,looseness=5] node[fmlabel,below] {$p/\mathsf{rst}(x_{r2})$} (s3);
  \draw[fmedge] (s3) -- node[fmlabel,right,pos=.4] {$q$} (s5);
  \draw[fmedge] (s3) -- node[fmlabel,below,sloped] {$x_{r2}>3,\;q$} (s2);
  % location 5
  \draw[fmedge] (s5) edge[loop above,looseness=5] node[fmlabel] {$p,\;q$} (s5);
  \draw[fmedge] (s5) -- node[fmlabel,right] {$x_{r2}>3,\;p,\;q,\;\mathit{other}$} (s2);
  \draw[fmedge] (s2) edge[loop right,looseness=5] node[fmlabel,right] {$p,\;q,\;\mathit{other}$} (s2);
\end{tikzpicture}
}
\caption{Optimized automaton for
$\Diamond_{\le5}((\Box_{\le3}(p\lor q))\,\mathsf U_{\le2}\,q)$ before
disjunctive-invariant elimination. Clocks $x_{u0}$, $x_{u1}$, and $x_{r2}$
belong to the three distinct primitive timed subformulas. The diagram is
after the initial \texttt{\_init\_} transition; guards implied by source
invariants are omitted. A comma-separated event list denotes separate
transitions. The $q$-edges from locations~4 and~3 with $x_{r2}>3$ are
unsatisfiable for the reasons given in the text.}
\label{fig:ex2-disj}
\end{figure*}

\subsubsection{Optimization and export of the F11 automaton}
The F11 automaton for
$ (\Diamond_{\le5}p)\land(\Box_{<2}\neg p) $ illustrates the interaction
between the generic passes. Spot returns five states and eleven transitions.
After one optimization round, contradictory action cubes have been removed,
making location~2 unreachable; dead resets are removed, and clocks $x$ and
$y$ are merged because every transition resets them synchronously. The
second round removes the now-dead resets of $y$. The final UPPAAL-oriented
automaton has four locations, seven transitions, one clock, and two location
invariants. Figure~\ref{fig:running-pipeline} shows the state after the first
round and the exported form. This example concerns the optimizer and export;
its temporal translation is handled separately.

\begin{figure*}[!t]
\centering
\resizebox{\textwidth}{!}{\begin{tikzpicture}[x=1cm,y=1cm,lab/.style={fmlabel}]
\begin{scope}[xshift=0cm]
  \node[fmnote,font=\small\bfseries] at (2.2,5.3) {After one optimization round};
  \node[fmstate] (l0) at (0.0,4.3) {$0$};
  \node[fmstate] (l1) at (0.0,2.3) {$1$};
  \node[fmstate] (l3) at (4.1,2.3) {$3$};
  \node[fmaccept] (l4) at (0.0,0.0) {$4$};
  \draw[fmedge] (-1.0,4.3) -- (l0);
  \draw[fmedge] (l0) -- node[lab,right] {$\neg p\;/\;\mathsf{rst}(x),\mathsf{rst}(y)$} (l1);
  \draw[fmedge] (l1) edge[loop left,looseness=5] node[lab,left] {$x\le5,\;\neg p$} (l1);
  \draw[fmedge] (l1) -- node[lab,above] {$2\le x\le5,\;\neg p$} (l3);
  \draw[fmedge] (l1) -- node[lab,left] {$2\le x\le5,\;p$} (l4);
  \draw[fmedge] (l3) edge[loop above,looseness=5] node[lab,above] {$x\le5,\;\neg p$} (l3);
  \draw[fmedge] (l3) -- node[lab,pos=.55,below right] {$x\le5,\;p$} (l4);
  \draw[fmedge] (l4) edge[loop right] node[lab,right] {$\top$} (l4);
\end{scope}

\begin{scope}[xshift=8.5cm]
  \node[fmnote,font=\small\bfseries] at (2.2,5.3) {Final TBA};
  \node[fmstate] (r0) at (0.0,4.3) {$0$};
  \node[fmstate,minimum size=11mm] (r1) at (0.0,2.3) {$\begin{array}{c}1\\[-1mm]x\le5\end{array}$};
  \node[fmstate,minimum size=11mm] (r3) at (4.1,2.3) {$\begin{array}{c}3\\[-1mm]x\le5\end{array}$};
  \node[fmaccept] (r4) at (0.0,0.0) {$4$};
  \draw[fmedge] (-1.0,4.3) -- (r0);
  \draw[fmedge] (r0) -- node[lab,right] {$\neg p\;/\;\mathsf{rst}(x)$} (r1);
  \draw[fmedge] (r1) edge[loop left,looseness=5] node[lab,left] {$x\le5,\;\neg p$} (r1);
  \draw[fmedge] (r1) -- node[lab,above] {$2\le x\le5,\;\neg p$} (r3);
  \draw[fmedge] (r1) -- node[lab,left] {$2\le x\le5,\;p$} (r4);
  \draw[fmedge] (r3) edge[loop above,looseness=5] node[lab,above] {$x\le5,\;\neg p$} (r3);
  \draw[fmedge] (r3) -- node[lab,pos=.55,below right] {$x\le5,\;p$} (r4);
  \draw[fmedge] (r4) edge[loop right] node[lab,right] {$\top$} (r4);
\end{scope}
\end{tikzpicture}
}
\caption{F11 after one optimization round and after UPPAAL-oriented export.
The first panel shows the removed unreachable location and the two
synchronously reset clocks; the final panel shows the merged clock and
synthesized invariants. Edges use the form ``guard, action / resets,'' and the
double circle marks the accepting location.}
\label{fig:running-pipeline}
\end{figure*}
